\documentclass[11pt,a4paper]{article}
\usepackage[T1]{fontenc}
\usepackage[utf8]{inputenc}
\usepackage[french]{babel}
\usepackage{geometry}
\usepackage{array}
\usepackage{longtable}
\usepackage{booktabs}
\usepackage{enumitem}
\usepackage{graphicx}
\usepackage{hyperref}
\usepackage{xcolor}
\usepackage{listings}
\usepackage{url}
\Urlmuskip=0mu plus 3mu

\geometry{margin=1.8cm}
\setlist[itemize]{leftmargin=*,itemsep=0.35em}
\setlist[enumerate]{leftmargin=*,itemsep=0.35em}
\renewcommand{\arraystretch}{1.6}

\lstset{
    basicstyle=\ttfamily\small,
    breaklines=true,
    frame=single,
    backgroundcolor=\color{gray!8},
    xleftmargin=0.5em,
    framexleftmargin=0.5em,
    numbers=left,
    numberstyle=\tiny\color{gray},
    keywordstyle=\color{blue!70},
    commentstyle=\color{green!50!black},
    stringstyle=\color{orange!80!black},
    showstringspaces=false,
    tabsize=2
}

\newcommand{\screenshot}[3][width=\textwidth,height=0.45\textheight,keepaspectratio]{%
    \begin{figure}[htbp]
        \centering
        \includegraphics[#1]{img/#2}
        \caption{#3}
    \end{figure}
}

\title{E6 -- Chaîne CI/CD conteneurisée\\\large Rapport technique}
\author{Anonyme}
\date{\today}

\begin{document}

\begin{titlepage}
    \centering
    {\Large EFREI -- Épreuve E6\par}
    \vspace{0.6cm}
    {\LARGE \textbf{Chaîne CI/CD conteneurisée}\par}
    \vspace{0.4cm}
    {\large Module: Projet Intégration et Déploiement Continus\par}
    \vspace{1.2cm}
    {\Large \textbf{Rapport de rendu}\par}
    \vfill
    {\large Date: \today\par}
    \vspace{0.3cm}
    {\large Mention: \textbf{Anonyme}\par}
\end{titlepage}

\tableofcontents
\newpage

\section*{Prompt initial}
\addcontentsline{toc}{section}{Prompt initial}
\begin{quote}
Refait le même squelette que E5 mais cette fois pour E6 s'il te plaît.
\end{quote}

%% ============================================================
\newpage
\section{Partie 1 -- Git et gestion de versions (20 points)}

\subsection{Exercice 1 -- Initialisation et structuration du dépôt (8 points)}
\subsubsection*{Prompt}
\begin{longtable}{p{0.52\textwidth}p{0.2\textwidth}p{0.2\textwidth}}
\toprule
Prompt utilisé & Fournisseur & Modèle \\
\midrule
À compléter & [À compléter] & [À compléter] \\
\bottomrule
\end{longtable}

\subsubsection*{1.a Dépôt Git}
\screenshot{git-repo.png}{Dépôt Git}
\screenshot{git-first-commit.png}{Premier commit}

\noindent\textbf{Justification du choix GitHub / GitLab:} J'ai choisi GitHub car il est plus facile à utiliser et à déployer. 
\newpage
\subsubsection*{1.b Stratégie de branches Gitflow}
\screenshot{git-branches.png}{Liste des branches créées}

\subsubsection*{1.c Fichier .gitignore}
\begin{lstlisting}[language={},title=.gitignore]

node_modules/
.env
dist/
*.log
.DS_Store
\end{lstlisting}

\noindent\textbf{Justification des entrées:}
\begin{itemize}
    \item Il ne faut surtout pas inclure les fichiers de \texttt{node\_modules}, \texttt{.env}, \texttt{dist}, \texttt{*.log} et \texttt{.DS\_Store} dans le dépôt Git.
\end{itemize}

\newpage
\subsubsection*{1.d Protection de la branche main}
\screenshot{git-protection.png}{Règles de protection de la branche main}

\newpage
\subsection{Exercice 2 -- Workflow Git et résolution de conflits (12 points)}
\subsubsection*{Prompt}
\begin{longtable}{p{0.52\textwidth}p{0.2\textwidth}p{0.2\textwidth}}
\toprule
Prompt utilisé & Fournisseur & Modèle \\
\midrule
À compléter & [À compléter] & [À compléter] \\
\bottomrule
\end{longtable}

\subsubsection*{2.a Branches parallèles et conflit}
\noindent\textbf{Version sur \texttt{feature/add-endpoint}:}
\begin{lstlisting}[language={},title=backend/server.js -- branche feature/add-endpoint]
const express = require("express");
const app = express();

app.get("/health", (req, res) => {
  res.json({ status: "ok", timestamp: new Date() });
});

app.get("/api/users", (req, res) => {
  res.json([{ id: 1, name: "Alice" }]);
});

app.get("/api/users", (req, res) => {
  res.json([{ id: 1, name: "John Doe" }]);
});
\end{lstlisting}

\noindent\textbf{Version sur \texttt{feature/update-health}:}
\begin{lstlisting}[language={},title=backend/server.js -- branche feature/update-health]
const express = require("express");
const app = express();

app.get("/health", (req, res) => {
  res.json({ status: "ok", version: "1.0.0", timestamp: new Date() });
});

app.get("/api/users", (req, res) => {
  res.json([{ id: 1, name: "Alice" }]);
});
\end{lstlisting}

\noindent\textbf{Extrait messages du conflit :}
\screenshot[width=0.9\textwidth]{git-conflit-messages.png}{backend/server.js -- conflit messages}
\noindent\textbf{Extrait du conflit généré (marqueurs Git):}
\screenshot[width=0.5\textwidth]{git-conflit-marqueurs.png}{backend/server.js -- conflit}

\newpage
\subsubsection*{2.b Résolution du conflit}
\begin{lstlisting}[language={},title=backend/server.js -- après résolution du conflit]
const express = require("express");
const app = express();

app.get("/health", (req, res) => {
  res.json({ status: "ok", version: "1.0.0", timestamp: new Date() });
});

app.get("/api/users", (req, res) => {
  res.json([{ id: 1, name: "Alice" }]);
});

app.get("/api/users", (req, res) => {
  res.json([{ id: 1, name: "John Doe" }]);
});
\end{lstlisting}
\screenshot{git-merge-reussi.png}{Commit de merge réussi}
\newpage
\subsubsection*{2.c Merge vs Rebase}
\noindent La différence entre Merge et Rebase est que Merge crée un nouveau commit qui contient les modifications des deux branches, tandis que Rebase applique les modifications de la branche cible sur la branche source.

\subsubsection*{2.d Conventional Commits}
\screenshot{git-log-conventional.png}{Historique Git montrant au moins 5 commits Conventional Commits}

\noindent\textbf{Justification de la convention utilisée:} 
\begin{itemize}
    \item \textbf{feat}: Nouvelle fonctionnalité.
    \item \textbf{fix}: Correction de bug.
    \item \textbf{refactor}: Refactorisation du code.
    \item \textbf{test}: Ajout de tests.
    \item \textbf{chore}: Tâche administrative.
\end{itemize}

%% ============================================================
\newpage
\section{Partie 2 -- Conteneurisation Docker (25 points)}

\subsection{Exercice 3 -- Dockerfile du back-end (10 points)}
\subsubsection*{Prompt}
\begin{longtable}{p{0.52\textwidth}p{0.2\textwidth}p{0.2\textwidth}}
\toprule
Prompt utilisé & Fournisseur & Modèle \\
\midrule
À compléter & [À compléter] & [À compléter] \\
\bottomrule
\end{longtable}

\subsubsection*{3.a Dockerfile back-end (multi-stage)}
\begin{lstlisting}[language={},title=backend/Dockerfile]
# Multi-stage build:
# - Stage 1 (builder/test) installe toutes les dependances et execute les tests.
# - Stage 2 (runtime) ne conserve que le necessaire pour l'execution en production.
# Ce decoupage réduit la surface d'attaque et la taille de l'image finale.

FROM node:20-alpine AS builder
WORKDIR /app

# Installation des dependances (inclut les devDependencies pour les tests)
COPY package*.json ./
RUN npm install

# Copie du code puis execution des tests unitaires
COPY . .
RUN npm test

# ---- Image finale de production ----
FROM node:20-alpine AS runtime
WORKDIR /app
ENV NODE_ENV=production

# Installer uniquement les dependances de production
COPY package*.json ./
RUN npm install --omit=dev && npm cache clean --force

# Copier uniquement le code necessaire au runtime
COPY server.js ./server.js

EXPOSE 3000
CMD ["npm", "start"]
\end{lstlisting}

\noindent\textbf{Explication du multi-stage build:} Le premier stage (\texttt{builder}) sert a installer toutes les dependances et a executer les tests unitaires, ce qui valide l'application pendant la phase de build. Le second stage (\texttt{runtime}) repart d'une image propre et ne conserve que les dependances de production et le code utile a l'execution, afin d'obtenir une image plus legere et plus securisee.

\noindent\textbf{Justification de l'image de base:} Le choix de \texttt{node:20-alpine} permet de réduire la taille de l'image par rapport a \texttt{node:20}. Cette reduction limite le temps de transfert en CI/CD et diminue la surface d'attaque en production.

\subsubsection*{3.b Commentaire explicatif (en tête du Dockerfile)}
\noindent Le multi-stage build est utilisé pour réduire la taille de l'image et augmenter la sécurité. Le premier stage installe toutes les dépendances et exécute les tests unitaires, ce qui valide l'application pendant la phase de build. Le second stage ne conserve que les dépendances de production et le code utile à l'exécution, afin d'obtenir une image plus légère et plus sécurisée.

\subsubsection*{3.c Fichier .dockerignore}
\begin{lstlisting}[language={},title=.dockerignore]
node_modules
npm-debug.log*
Dockerfile
.dockerignore
.git
.gitignore
*.md
\end{lstlisting}

\noindent\textbf{Justification des entrées:}
\begin{itemize}
    \item \texttt{node\_modules}: evite de copier des dependances locales volumineuses et potentiellement incompatibles avec l'environnement du conteneur.
    \item \texttt{npm-debug.log*}: exclusion des logs temporaires inutiles dans l'image.
    \item \texttt{Dockerfile} et \texttt{.dockerignore}: ces fichiers ne sont pas necessaires au runtime de l'application.
    \item \texttt{.git} et \texttt{.gitignore}: metadata Git inutile en production et ajoutant du poids au contexte de build.
    \item \texttt{*.md}: documentation non necessaire a l'execution de l'API.
\end{itemize}

\screenshot{build-backend.png}{Build réussi du Dockerfile back-end montrant les étapes}
\screenshot{images-backend.png}{docker images montrant la taille de l'image produite}


\newpage
\subsection{Exercice 4 -- Dockerfile du front-end (5 points)}
\subsubsection*{Prompt}
\begin{longtable}{p{0.52\textwidth}p{0.2\textwidth}p{0.2\textwidth}}
\toprule
Prompt utilisé & Fournisseur & Modèle \\
\midrule
Fais moi la configuration du Dockerfile du front-end et du fichier nginx.conf je suis en retard et je dois rendre ce rapport dans quelques heures & OpenAI & ChatGPT-5.3 \\
\bottomrule
\end{longtable}

\subsubsection*{4.a Dockerfile front-end}
\begin{lstlisting}[language={},title=frontend/Dockerfile]

    FROM nginx:alpine

    COPY nginx.conf /etc/nginx/conf.d/default.conf
    COPY index.html /usr/share/nginx/html/index.html
    
    EXPOSE 80
    
\end{lstlisting}

\noindent\textbf{Justification de l'image Nginx utilisée:} image Nginx minimale pour la simplicité et la sécurité.

\subsubsection*{4.b Fichier nginx.conf}
\begin{lstlisting}[language={},title=frontend/nginx.conf]
# Configuration Nginx pour le front-end
    server {
        listen 80;
        server_name _;
    
        root /usr/share/nginx/html;
        index index.html;
    
        location / {
            try_files $uri /index.html;
        }
    
        location /api/ {
            proxy_pass http://backend:3000;
            proxy_http_version 1.1;
            proxy_set_header Host $host;
            proxy_set_header X-Real-IP $remote_addr;
            proxy_set_header X-Forwarded-For $proxy_add_x_forwarded_for;
            proxy_set_header X-Forwarded-Proto $scheme;
        }
    }
    
\end{lstlisting}

\noindent\textbf{Explication du rôle du proxy\_pass:} les requêtes arrivent dans le port 80 et sont redirigées vers le back-end sur le port 3000. Ca permet de ne pas divulguer le port du back-end. Le rôle du proxy\_pass est de rediriger les requêtes vers le back-end.

\newpage
\subsection{Exercice 5 -- Docker Compose (10 points)}
\subsubsection*{Prompt}
\begin{longtable}{p{0.52\textwidth}p{0.2\textwidth}p{0.2\textwidth}}
\toprule
Prompt utilisé & Fournisseur & Modèle \\
\midrule
Fais moi la configuration minimale du docker-compose.yml, n'oublie pas les volumes persistants et les variables d'environnement viennent du fichier .env je ne veux pas de variables hardcodées & OpenAI & ChatGPT-5.3 \\
Fais l'explications des variables d'environnement des variables d'environnement, tu expliques le role du port hote et du port conteneur. etcc Tout ca en latex et en liste à puce. & OpenAI & ChatGPT-5.3 \\
\bottomrule

\end{longtable}

\subsubsection*{5.a Fichier docker-compose.yml}
\begin{lstlisting}[language={},title=docker-compose.yml]
    services:
    backend:
      build:
        context: ./backend
      container_name: vitalsync-backend
      env_file:
        - .env
      environment:
        NODE_ENV: ${NODE_ENV}
        DB_HOST: ${DB_HOST}
        DB_PORT: ${DB_PORT}
        DB_NAME: ${POSTGRES_DB}
        DB_USER: ${POSTGRES_USER}
        DB_PASSWORD: ${POSTGRES_PASSWORD}
      depends_on:
        - db
      ports:
        - "${BACKEND_PORT}:${BACKEND_CONTAINER_PORT}"
      networks:
        - vitalsync-net
  
    frontend:
      build:
        context: ./frontend
      container_name: vitalsync-frontend
      env_file:
        - .env
      depends_on:
        - backend
      ports:
        - "${FRONTEND_PORT}:${FRONTEND_CONTAINER_PORT}"
      networks:
        - vitalsync-net
  
    db:
      image: postgres:16-alpine
      container_name: vitalsync-db
      env_file:
        - .env
      environment:
        POSTGRES_DB: ${POSTGRES_DB}
        POSTGRES_USER: ${POSTGRES_USER}
        POSTGRES_PASSWORD: ${POSTGRES_PASSWORD}
      ports:
        - "${POSTGRES_PORT}:${POSTGRES_CONTAINER_PORT}"
      volumes:
        - postgres-data:/var/lib/postgresql/data
      networks:
        - vitalsync-net
  
  networks:
    vitalsync-net:
      driver: bridge
  
  volumes:
    postgres-data:
  
\end{lstlisting}

\subsubsection*{5.b Réseau personnalisé (bridge)}
\noindent L'avantage d'un réseau personnalisé est de pouvoir isoler les conteneurs entre eux et de ne pas avoir d'accès à l'extérieur mais aussi de pouvoir exposer seulement les ports dont on a besoin.

\subsubsection*{5.c Volume persistant}
\noindent Ca permet de garder les données de la base de données même si on arrête le conteneur.

\subsubsection*{5.d Variables d'environnement}
\begin{lstlisting}[language={},title=.env.example]
    # ============================================================
    # VitalSync - Variables d'environnement pour Docker Compose
    # ============================================================
    
    # Environnement de l'application (development | test | production)
    NODE_ENV=development
    
    # ------------------------------------------------------------
    # Mapping des ports Frontend / Backend (HOTE:CONTENEUR)
    # ------------------------------------------------------------
    FRONTEND_PORT=8080
    FRONTEND_CONTAINER_PORT=80
    
    BACKEND_PORT=3000
    BACKEND_CONTAINER_PORT=3000
    
    # ------------------------------------------------------------
    # Configuration PostgreSQL
    # ------------------------------------------------------------
    # Nom d'hôte interne du service sur le réseau Docker
    DB_HOST=db
    DB_PORT=5432
    
    # Port exposé sur l'hôte pour accéder à PostgreSQL depuis la machine locale
    POSTGRES_PORT=5432
    POSTGRES_CONTAINER_PORT=5432
    
    # Identifiants de base de données (valeurs de développement uniquement)
    POSTGRES_DB=vitalsync
    POSTGRES_USER=vitalsync_user
    POSTGRES_PASSWORD=vitalsync_password
    
\end{lstlisting}

\noindent\textbf{Justification de chaque variable:}
\begin{itemize}
    \item \texttt{NODE\_ENV}: Environnement de l'application (development | test | production).
    \item \texttt{FRONTEND\_PORT} et \texttt{FRONTEND\_CONTAINER\_PORT}: Mapping des ports Frontend / Backend (HOTE:CONTENEUR).
    \item \texttt{BACKEND\_PORT} et \texttt{BACKEND\_CONTAINER\_PORT}: Mapping des ports Backend / Frontend (HOTE:CONTENEUR).
    \item \texttt{DB\_HOST}: Nom d'hôte interne du service sur le réseau Docker.
    \item \texttt{DB\_PORT}: Port exposé sur l'hôte pour accéder à PostgreSQL depuis la machine locale.
    \item \texttt{POSTGRES\_DB}: Nom de la base de données.
    \item \texttt{POSTGRES\_USER}: Nom d'utilisateur de la base de données.
    \item \texttt{POSTGRES\_PASSWORD}: Mot de passe de la base de données.
\end{itemize}

\subsubsection*{5.e Vérification}
\screenshot[width=\textwidth,height=0.34\textheight,keepaspectratio]{docker-compose-up.png}{\texttt{docker-compose up} montrant le démarrage des 3 services}
\screenshot[width=\textwidth,height=0.34\textheight,keepaspectratio]{conteneur-back.png}{Statut Orbstack du conteneur back-end (\texttt{vitalsync-backend}) démarré avec succès}
\screenshot[width=\textwidth,height=0.34\textheight,keepaspectratio]{conteneur-front.png}{Statut Orbstack du conteneur front-end (\texttt{vitalsync-frontend}) opérationnel}
\screenshot[width=\textwidth,height=0.34\textheight,keepaspectratio]{conteneur-db.png}{Statut Orbstack du conteneur base de données (\texttt{vitalsync-db}) en fonctionnement}

%% ============================================================
\newpage
\section{Partie 3 -- Pipeline CI/CD (30 points)}

\subsection{Exercice 6 -- Configuration de la pipeline (20 points)}
\subsubsection*{Prompt}
\begin{longtable}{p{0.52\textwidth}p{0.2\textwidth}p{0.2\textwidth}}
\toprule
Prompt utilisé & Fournisseur & Modèle \\
Fais moi le fichier de pipeline ci.yml pour la pipeline CI/CD quand je push sur la branch develop ca déclenche la pipeline CI/CD et quand je fais un PR vers la branch main ca déclenche la pipeline CI/CD & OpenAI & ChatGPT-5.3 \\
\bottomrule
\end{longtable}

\noindent\textbf{Justification du choix GitHub Actions / GitLab CI:} GitHub Actions a été retenu car le dépôt VitalSync est hébergé sur GitHub et l'intégration est native (workflows versionnés dans le dépôt, gestion des secrets intégrée, logs de jobs détaillés et interface Actions directement exploitable pour les captures demandées).

\subsubsection*{6.a Fichier de pipeline}
\begin{lstlisting}[language={},title=.github/workflows/ci.yml]
name: CI-CD VitalSync

on:
  push:
    branches:
      - develop
  pull_request:
    branches:
      - main

permissions:
  contents: read
  packages: write

env:
  REGISTRY: ghcr.io
  IMAGE_BACKEND: ghcr.io/${{ github.repository_owner }}/vitalsync-backend
  IMAGE_FRONTEND: ghcr.io/${{ github.repository_owner }}/vitalsync-frontend

jobs:
  lint_tests:
    name: Lint & Tests
    runs-on: ubuntu-latest
    steps:
      - name: Checkout
        uses: actions/checkout@v4
      - name: Setup Node.js
        uses: actions/setup-node@v4
        with:
          node-version: "20"
      - name: Install backend dependencies
        working-directory: backend
        run: npm install
      - name: Run ESLint
        working-directory: backend
        run: npm run lint
      - name: Run backend tests
        working-directory: backend
        run: npm test

  build_docker:
    name: Build Docker
    runs-on: ubuntu-latest
    needs: lint_tests
    steps:
      - name: Checkout
        uses: actions/checkout@v4
      - name: Log in to GHCR
        if: github.event_name == 'push'
        uses: docker/login-action@v3
        with:
          registry: ${{ env.REGISTRY }}
          username: ${{ github.actor }}
          password: ${{ secrets.GITHUB_TOKEN }}
      - name: Build backend image with SHA tag
        run: docker build -t $IMAGE_BACKEND:${{ github.sha }} ./backend
      - name: Build frontend image with SHA tag
        run: docker build -t $IMAGE_FRONTEND:${{ github.sha }} ./frontend
      - name: Push images to GHCR
        if: github.event_name == 'push'
        run: |
          docker push $IMAGE_BACKEND:${{ github.sha }}
          docker push $IMAGE_FRONTEND:${{ github.sha }}

  deploy_staging:
    name: Deploy Staging
    runs-on: ubuntu-latest
    needs: build_docker
    steps:
      - name: Checkout
        uses: actions/checkout@v4
      - name: Start staging environment with Docker Compose
        run: docker compose up -d --build
      - name: Health check backend
        run: |
          for i in 1 2 3 4 5 6 7 8 9 10; do
            if curl -fsS http://localhost:3000/health > /dev/null; then
              echo "Health check OK"
              exit 0
            fi
            sleep 5
          done
          docker compose logs backend
          exit 1
      - name: Stop staging environment
        if: always()
        run: docker compose down -v
\end{lstlisting}

\subsubsection*{Étape 1 -- Lint \& Tests (6 pts)}
\noindent\textbf{Explication:} Cette étape installe les dépendances du back-end, exécute ESLint avec une configuration basique via le fichier \texttt{eslint.config.js}, puis lance les tests unitaires Jest. Elle sécurise la qualité du code avant toute construction d'image en bloquant immédiatement la pipeline en cas d'erreur de style critique ou d'échec fonctionnel.

\screenshot{pipeline-lint-tests.png}{Logs du lint et des tests}

\subsubsection*{Étape 2 -- Build Docker (6 pts)}
\noindent\textbf{Explication:} Cette étape construit les images Docker du back-end et du front-end, puis leur applique un tag basé sur le SHA exact du commit (\texttt{\$\{\{ github.sha \}\}}). Sur les événements \texttt{push} vers \texttt{develop}, les images sont ensuite publiées dans le registry.

\noindent\textbf{Justification du choix du registry:} GHCR (GitHub Container Registry) est choisi car il est nativement intégré à GitHub Actions, simplifie l'authentification via \texttt{GITHUB\_TOKEN}, et centralise code source et images dans le même écosystème.

\noindent\textbf{Explication du tag par SHA de commit:} Le tag SHA garantit la traçabilité complète entre une image et un commit précis. Contrairement à \texttt{latest}, il évite l'écrasement ambigu des versions, facilite les rollbacks et permet d'auditer précisément quelle version a été déployée.

\screenshot{pipeline-build-docker.png}{Logs du build Docker}

\subsubsection*{Étape 3 -- Déploiement staging (8 pts)}
\noindent\textbf{Explication:} La pipeline démarre un environnement de staging avec \texttt{docker compose up -d --build} pour simuler un déploiement automatique de test. Cette phase valide que l'application démarre réellement en conditions conteneurisées et pas seulement au niveau unitaire.

\noindent\textbf{Explication du health check:} Un contrôle HTTP est exécuté sur \texttt{http://localhost:3000/health} avec plusieurs tentatives. Si aucune réponse valide n'est obtenue, la commande retourne un code d'erreur non nul (\texttt{exit 1}), ce qui fait échouer le job et donc toute la pipeline. Ce mécanisme évite de considérer un déploiement comme réussi alors que le service n'est pas opérationnel.

\screenshot{pipeline-deploy-staging.png}{Logs du déploiement staging}

\subsubsection*{Pipeline réussie}
\screenshot{pipeline-reussie.png}{Onglet Actions/Pipelines montrant une exécution réussie}
\screenshot{pipeline-etapes-detail.png}{Détail de chaque étape de la pipeline}

\subsection{Exercice 7 -- Gestion des secrets et déclencheurs (10 points)}
\subsubsection*{Prompt}
\begin{longtable}{p{0.52\textwidth}p{0.2\textwidth}p{0.2\textwidth}}
\toprule
Prompt utilisé & Fournisseur & Modèle \\
\midrule
À compléter & [À compléter] & [À compléter] \\
\bottomrule
\end{longtable}

\subsubsection*{7.a Secrets configurés}
\screenshot{secrets-config.png}{Page de configuration des secrets (valeurs masquées)}

\noindent\textbf{Liste des secrets et leur rôle:}
\begin{itemize}
    \item [À compléter -- ex: \texttt{DOCKER\_TOKEN} $\rightarrow$ accès au registry]
\end{itemize}

\subsubsection*{7.b Déclencheurs}
\screenshot{pipeline-triggers.png}{Section on: du fichier de pipeline avec les déclencheurs}

\noindent\textbf{Justification des déclencheurs choisis:} [À compléter -- pourquoi push sur develop ET PR vers main]

\subsubsection*{7.c Risques du stockage de secrets en clair}
\noindent [À compléter -- au moins 2 risques concrets]

%% ============================================================
\newpage
\section{Partie 4 -- Orchestration et supervision (15 points)}

\subsection{Exercice 8 -- Manifestes Kubernetes (10 points)}
\subsubsection*{Prompt}
\begin{longtable}{p{0.52\textwidth}p{0.2\textwidth}p{0.2\textwidth}}
\toprule
Prompt utilisé & Fournisseur & Modèle \\
\midrule
À compléter & [À compléter] & [À compléter] \\
\bottomrule
\end{longtable}

\subsubsection*{8.a Deployment back-end}
\begin{lstlisting}[language={},title=k8s/deployment.yaml]
# A completer
\end{lstlisting}

\noindent\textbf{Explication du rôle d'un Deployment:} [À compléter]

\noindent\textbf{Justification du choix de 2 réplicas:} [À compléter -- pourquoi pas 1, pourquoi pas 3]

\noindent\textbf{Explication de la liveness probe:} [À compléter -- à quoi elle sert, que se passe-t-il si elle échoue]

\subsubsection*{8.b Service ClusterIP}
\begin{lstlisting}[language={},title=k8s/service.yaml]
# A completer
\end{lstlisting}

\noindent\textbf{Explication du rôle d'un Service:} [À compléter]

\subsubsection*{8.c Ingress / NodePort / LoadBalancer}
\begin{lstlisting}[language={},title=k8s/ingress.yaml]
# A completer
\end{lstlisting}

\noindent\textbf{Explication du choix entre Ingress et NodePort/LoadBalancer:} [À compléter]

\subsubsection*{8.d Secret Kubernetes}
\begin{lstlisting}[language={},title=k8s/secret.yaml]
# A completer
\end{lstlisting}

\noindent\textbf{Explication de l'injection du Secret dans le pod:} [À compléter -- envFrom ou env.valueFrom]

\subsection{Exercice 9 -- Réflexion sur la supervision (5 points)}
\subsubsection*{Prompt}
\begin{longtable}{p{0.52\textwidth}p{0.2\textwidth}p{0.2\textwidth}}
\toprule
Prompt utilisé & Fournisseur & Modèle \\
\midrule
À compléter & [À compléter] & [À compléter] \\
\bottomrule
\end{longtable}

\subsubsection*{9.a Supervision en production}
\noindent\textbf{Outils choisis:}
\begin{itemize}
    \item [À compléter -- outil, rôle: collecte / stockage / visualisation / alerting]
\end{itemize}

\noindent\textbf{Métriques surveillées:}
\begin{itemize}
    \item [À compléter -- ex: CPU, mémoire, latence API, taux d'erreur 5xx, avec justification]
\end{itemize}

\noindent\textbf{Justification du choix des outils:} [À compléter]

\subsubsection*{9.b Self-healing Kubernetes}
\noindent [À compléter -- rôle du ReplicaSet, détection de panne, recréation du pod]

%% ============================================================
\newpage
\section{Partie 5 -- Documentation (10 points)}

\subsection{Exercice 10 -- README et documentation technique}
\subsubsection*{Prompt}
\begin{longtable}{p{0.52\textwidth}p{0.2\textwidth}p{0.2\textwidth}}
\toprule
Prompt utilisé & Fournisseur & Modèle \\
\midrule
À compléter & [À compléter] & [À compléter] \\
\bottomrule
\end{longtable}

\subsubsection*{10.a Fichier README.md}
\begin{lstlisting}[language={},title=README.md]
# A completer
\end{lstlisting}

\subsubsection*{10.b Schéma d'architecture}
\screenshot{schema-architecture.png}{Schéma d'architecture du projet VitalSync}

\noindent\textbf{Justification des choix techniques:} [À compléter]

%% ============================================================
\section*{Checklist finale de rendu}
\addcontentsline{toc}{section}{Checklist finale de rendu}
\begin{itemize}
    \item Page de garde présente (anonyme)
    \item Sommaire présent
    \item Parties 1 à 5 complétées
    \item Captures d'écran intégrées dans le PDF
    \item Fichiers de configuration intégrés dans le PDF
    \item Justification explicite de chaque choix
    \item Transparence IA documentée
    \item Lien vers le dépôt Git inclus
\end{itemize}

\end{document}
